\documentclass[11pt,a4paper]{article}

% ---------------------------------------------------------------------------
% Packages
% ---------------------------------------------------------------------------
\usepackage[utf8]{inputenc}
\usepackage[T1]{fontenc}
\usepackage[margin=1in]{geometry}
\usepackage{graphicx}
\usepackage{xcolor}
\usepackage{booktabs}
\usepackage{listings}
\usepackage{float}
\usepackage{enumitem}
\usepackage{caption}
\usepackage[hidelinks]{hyperref}

% ---------------------------------------------------------------------------
% A screenshot helper: includes the image if it exists, otherwise draws a
% labelled placeholder box so the report always compiles. Replace the files in
% Report/figures/ to fill the report in.
% ---------------------------------------------------------------------------
\newcommand{\screenshot}[3]{%
  \begin{figure}[H]
    \centering
    \IfFileExists{figures/#1}{%
      \includegraphics[width=0.85\linewidth]{figures/#1}%
    }{%
      \fbox{\parbox[c][4.5cm][c]{0.85\linewidth}{\centering\ttfamily
        [ screenshot placeholder ]\\[2pt] \detokenize{figures/#1} \\[6pt]
        \normalfont\itshape (drop your snapshot here)}}%
    }
    \caption{#2}
    \label{fig:#3}
  \end{figure}%
}

% Listing style for commands / code
\definecolor{codebg}{gray}{0.96}
\lstset{
  basicstyle=\ttfamily\small,
  backgroundcolor=\color{codebg},
  breaklines=true,
  frame=single,
  framesep=4pt,
  columns=fullflexible,
  keepspaces=true,
  showstringspaces=false,
}

% ---------------------------------------------------------------------------
% Title
% ---------------------------------------------------------------------------
\title{\textbf{TCP SYN-Flood Attack and Defense}\\[4pt]
  \large A Controlled Denial-of-Service Experiment and Two Countermeasures}
\author{Course: CSE 406 --- Computer Security Sessional\\
  Department of Computer Science and Engineering, BUET\\[4pt]
  Group members: \emph{[Iftekhar Mahi (2105003)]}, \emph{[teammate 2]}, \emph{[teammate 3]}}
\date{\today}

\begin{document}
\maketitle

\begin{abstract}
This report documents a controlled TCP SYN-flood (half-open connection) attack
carried out entirely between our own devices on an isolated wireless network,
and the two countermeasures we implemented against it. We describe the attack
tooling, the exact steps performed, the output observed on the attacker, the
victim (server), and the measurement machine, and we judge whether the attack
succeeded and why. We then present two defenses: an application-level per-source
detector-and-blocker (\texttt{server/defense.py}) whose main weakness is that it
can also block legitimate users, and a stateless SYN-cookie defense
(\texttt{server/syn\_cookies.py}, mirroring the Linux kernel's
\texttt{net.ipv4.tcp\_syncookies}) that mitigates the flood without such
collateral damage.
\end{abstract}

\tableofcontents
\bigskip

% ===========================================================================
\section{Objective and Ethical Scope}
% ===========================================================================
The goal was to reproduce, observe, and defend against a TCP SYN-flood
denial-of-service attack. A SYN flood exploits the TCP three-way handshake: the
attacker sends a stream of \texttt{SYN} packets, usually with spoofed source
addresses, and never completes the handshake. Each such packet forces the victim
to allocate a half-open (\texttt{SYN-RECEIVED}) connection and reply with a
\texttt{SYN-ACK}. When these half-open entries fill the listen backlog, the
server can no longer accept connections from legitimate clients.

\paragraph{Ethical scope.} All traffic was confined to hardware we own on a
private hotspot with no other devices attached. Source addresses were spoofed
only within an unused slice of our own subnet, and the attack was never left
running unattended. No third-party system was targeted at any point.

% ===========================================================================
\section{Experimental Setup}
% ===========================================================================
The experiment used two physical laptops on a single NetworkManager shared
Wi-Fi network (\texttt{10.42.0.0/24}). One machine plays the \emph{server}
(victim) and also hosts the ``legitimate user'' measurement probe; the other is
the \emph{attacker}.

\begin{table}[H]
\centering
\caption{Roles, addresses, and software.}
\begin{tabular}{@{}lll@{}}
\toprule
\textbf{Role} & \textbf{Address} & \textbf{Runs} \\ \midrule
Server / victim & \texttt{10.42.0.81} & nginx (port 80), \texttt{server/monitor.py}, defenses \\
Server as ``real user'' & \texttt{10.42.0.81} & \texttt{client/sustained\_http\_load.py} probe \\
Attacker & \texttt{10.42.0.157} & \texttt{client/generator.py} (raw sockets, root) \\
Spoofed sources & \texttt{10.42.0.193--254} & synthetic addresses in the SYN packets \\ \bottomrule
\end{tabular}
\end{table}

The spoofed-source pool (\texttt{10.42.0.192/26}) was chosen so that it excludes
every real device on the network --- the gateway (\texttt{.1}), the server
(\texttt{.81}), and the attacker (\texttt{.157}) --- so that a spoofed
\texttt{SYN-ACK} can never be answered with a \texttt{RST} by a real host, which
would otherwise prematurely free a half-open slot.

\paragraph{Attack tool.} \texttt{client/generator.py} builds complete 40-byte
IPv4/TCP \texttt{SYN} packets by hand (custom IPv4 and TCP headers with correct
Internet checksums) and sends them over a raw socket. It cycles every
non-privileged TCP source port (1024--65535) for one synthetic source IP before
advancing to the next, giving each half-open connection a distinct identity.

\paragraph{Measurement tools.}
\begin{itemize}[nosep]
  \item \texttt{server/monitor.py} --- samples the victim every 0.5\,s into a CSV:
    \texttt{nginx\_active}, \texttt{port\_80\_listening}, \texttt{syn\_recv\_count},
    \texttt{established\_count}, \texttt{load\_1m}, \texttt{mem\_available\_kib}.
  \item \texttt{client/sustained\_http\_load.py} --- a legitimate-user probe that
    repeatedly fetches \texttt{http://10.42.0.81/} and records per-request
    success, latency, and status.
  \item \texttt{ss}, \texttt{tcpdump}, and \texttt{nstat} for live TCP-state,
    packet-level, and kernel-counter evidence.
\end{itemize}

\screenshot{topology.png}{Network topology of the two-device experiment.}{topology}

% ===========================================================================
\section{(a) Steps of the Attack, with Snapshots}
% ===========================================================================
The attack and its measurement were performed in the following order.

\subsection*{Step 0 --- Reachability and isolation check}
From the attacker we confirmed the victim was reachable and that the access
point did not isolate clients:
\begin{lstlisting}[language=bash]
ping -c3 10.42.0.81
\end{lstlisting}

\subsection*{Step 1 --- Baseline (no attack)}
On the server we started the health sampler and, from the server's own hotspot
address, ran the legitimate-user probe to record normal behaviour:
\begin{lstlisting}[language=bash]
# server, terminal 1
python3 -m server.monitor --interval 0.5 --output results/baseline-monitor.csv
# server, terminal 2 (the "real user")
python3 -m client.sustained_http_load --url http://10.42.0.81/ \
        --duration 20 --timeout 2 --concurrency 10 --output results/baseline-http.csv
\end{lstlisting}

\subsection*{Step 2 --- Launch the SYN flood}
On the attacker we started the flood (beginning with a low rate to confirm
delivery, then scaling up):
\begin{lstlisting}[language=bash]
sudo python3 -m client.generator --continuous --rate 1000 --workers 8
\end{lstlisting}

\screenshot{attacker_generator.png}{Attacker terminal: \texttt{client/generator.py} sending spoofed SYNs and reporting the actual send rate.}{attacker}

\subsection*{Step 3 --- Observe the victim under attack}
While the flood ran, we watched half-open connections accumulate and re-ran the
legitimate-user probe:
\begin{lstlisting}[language=bash]
# server: live count of half-open connections on port 80
watch -n0.5 'ss -Hnt state syn-recv "( sport = :80 )" | wc -l'
# server: packet-level proof SYNs are arriving
sudo tcpdump -ni <iface> 'tcp port 80 and tcp[tcpflags] & tcp-syn != 0 and tcp[tcpflags] & tcp-ack == 0'
# server: the "real user" probe during the attack window
python3 -m client.sustained_http_load --url http://10.42.0.81/ \
        --duration 20 --timeout 2 --concurrency 10 --output results/attack-http.csv
\end{lstlisting}

\screenshot{server_tcpdump.png}{Victim: \texttt{tcpdump} showing a stream of incoming SYNs from spoofed \texttt{10.42.0.193--254} sources.}{tcpdump}
\screenshot{server_ss_synrecv.png}{Victim: the \texttt{SYN-RECEIVED} count climbing toward the listen backlog during the attack.}{ssrecv}

\subsection*{Step 4 --- Stop and collect evidence}
We stopped the generator (\texttt{Ctrl+C}) and the monitor, then compared the
baseline and attack CSVs and the kernel counters (\texttt{nstat -az} before/after).

% ===========================================================================
\section{(c) Observed Output on Each Machine}
% ===========================================================================
\subsection{Attacker PC}
The generator reported the number of SYN packets sent, the requested rate, the
elapsed time, and the \emph{actual} achieved rate (Wi-Fi and the shared-hotspot
bandwidth cap the sustained rate below the request). A representative line:
\begin{lstlisting}
Sending to 10.42.0.81:80 with 8 worker(s); Ctrl+C stops sending.
sent NNNNN SYN packets with synthetic sources from 10.42.0.192/26 to 10.42.0.81:80
requested rate: 1000 packets/second; elapsed: T seconds; actual rate: R packets/second
\end{lstlisting}

\subsection{Victim PC (server)}
On the victim the effect was visible three ways:
\begin{itemize}[nosep]
  \item \texttt{ss \dots state syn-recv} --- the half-open count rose from
    $\approx 0$ at rest toward the nginx listen backlog (511) and stayed there
    while the flood ran, then drained after it stopped.
  \item \texttt{tcpdump} --- a continuous stream of \texttt{Flags [S]} packets
    from many distinct \texttt{10.42.0.193--254} sources to \texttt{10.42.0.81:80}.
  \item \texttt{server/monitor.py} CSV --- the \texttt{syn\_recv\_count} column
    spiked during the attack window (plotted in Figure~\ref{fig:plot}).
\end{itemize}

\screenshot{monitor_csv_plot.png}{Victim: \texttt{syn\_recv\_count} over time from \texttt{server/monitor.py} (baseline vs.\ attack).}{plot}

\subsection{Measurement PC (the ``real user'' probe)}
The legitimate-user probe on the server's own address showed the
service-level impact: comparing \texttt{baseline-http.csv} with
\texttt{attack-http.csv}, the success rate and/or the request latency
(Table~\ref{tab:http}) is the ground truth for whether real users were affected.

\screenshot{victim_http_fail.png}{Legitimate-user probe degrading (timeouts / higher latency) during the unprotected flood.}{httpfail}

\begin{table}[H]
\centering
\caption{Legitimate HTTP probe: baseline vs.\ attack (fill from your CSVs).}
\label{tab:http}
\begin{tabular}{@{}lccc@{}}
\toprule
\textbf{Condition} & \textbf{Success rate} & \textbf{Median latency} & \textbf{Timeouts} \\ \midrule
Baseline (no attack)        & \_\_ / \_\_ & \_\_ ms & \_\_ \\
Unprotected (attack)        & \_\_ / \_\_ & \_\_ ms & \_\_ \\
Protected: SYN cookies      & \_\_ / \_\_ & \_\_ ms & \_\_ \\ \bottomrule
\end{tabular}
\end{table}

% ===========================================================================
\section{(b) Was the Attack Successful?}
% ===========================================================================
We treat the attack as \textbf{successful} only if two things hold together, so
that mere packet generation is not mistaken for impact:
\begin{enumerate}[nosep]
  \item \textbf{The victim held the attack state:} \texttt{syn\_recv\_count}
    spiked and stayed near the listen backlog during the attack window; and
  \item \textbf{There was a measurable service effect:} the legitimate probe's
    success rate fell and/or its latency rose (Table~\ref{tab:http}), \emph{or},
    if HTTP stayed healthy, the kernel's \texttt{TcpExtListenDrops} /
    \texttt{TcpExtSyncookiesSent} counters rose sharply --- showing the packets
    arrived and were dropped or absorbed rather than silently lost.
\end{enumerate}

\paragraph{Why it succeeded (when it did).} The half-open queue is a finite
kernel resource. Because every spoofed source never sends the final \texttt{ACK},
each \texttt{SYN} pins one backlog slot for the full \texttt{SYN-RECEIVED}
timeout; a fast enough arrival rate keeps the queue saturated, so genuine
handshakes are refused. Our use of many distinct source IP/port pairs made the
half-open entries distinct rather than collapsing onto one connection.

\paragraph{Why it can fail or under-perform.} Several real factors limited the
attack and are worth reporting honestly:
\begin{itemize}[nosep]
  \item \textbf{Wi-Fi / hotspot limits:} the shared radio caps the sustained
    packet rate well below the requested 1000/s, so the queue may not fully
    saturate.
  \item \textbf{Kernel SYN cookies:} if \texttt{net.ipv4.tcp\_syncookies} is
    enabled (the default on modern Linux), the victim answers backlog overflow
    with cookies and keeps serving --- the attack ``lands'' but does not deny
    service (this is exactly the protected case in \S\ref{sec:def}).
  \item \textbf{Access-point filtering:} some hotspots drop packets whose source
    IP does not match the associated client, silently discarding spoofed SYNs
    before they reach the victim.
\end{itemize}

% ===========================================================================
\section{(d) Countermeasures We Designed}
\label{sec:def}
% ===========================================================================
We implemented and evaluated two defenses with complementary trade-offs.

\subsection{Countermeasure 1 --- \texttt{server/defense.py} (SYN Guard): a per-source detector/blocker}
\texttt{defense.py} runs on the victim on a timer. Each tick it reads the
half-open connections from \texttt{ss}, counts them per source IP, and for any
source whose half-open count meets a threshold it inserts a temporary
\texttt{iptables} rule dropping further \texttt{SYN}s from that address, logging
every action to a CSV and lifting the block after a cooldown:
\begin{lstlisting}[language=bash]
sudo python3 -m server.defense --interval 0.5 --threshold 5 \
     --block-seconds 30 --output results/protected-defense-events.csv
\end{lstlisting}

\screenshot{defense_block.png}{\texttt{defense.py} detecting an over-threshold source and inserting an \texttt{iptables} DROP rule.}{block}

\paragraph{\textcolor{red}{Important limitation --- it also blocks legitimate
users.}} \texttt{defense.py} decides guilt purely by \emph{source IP address},
so its mitigation is inseparable from collateral damage:
\begin{itemize}[nosep]
  \item \textbf{It blocks innocent addresses.} A SYN flood uses \emph{spoofed}
    source IPs. When \texttt{defense.py} blocks a flagged address, it is blocking
    an address the attacker merely \emph{claimed} to be --- which may belong to a
    real, innocent host (or the whole subnet behind a NAT). That legitimate user
    is now firewalled off even though they did nothing.
  \item \textbf{It cannot cope with a widely distributed/spoofed flood.} If each
    fake source stays just under the threshold, or if thousands of distinct
    spoofed IPs are used, the defense either never triggers or must block huge
    ranges of addresses --- again sweeping up legitimate clients.
  \item \textbf{Threshold is a blunt instrument.} Set it low and you block bursty
    but legitimate clients; set it high and a real flood slips under it. There is
    no setting that cleanly separates attacker from victim, because the signal it
    uses (source IP) is exactly the field the attacker forges.
\end{itemize}
In short, \texttt{defense.py} is a useful, observable, and tunable detector
against a \emph{few real, high-volume} sources, but by construction it can punish
legitimate users, so it is not a safe standalone defense against a spoofed SYN
flood.

\subsection{Countermeasure 2 --- \texttt{server/syn\_cookies.py}: stateless SYN cookies}
The second defense removes the resource the attack targets instead of trying to
identify attackers. \texttt{syn\_cookies.py} reimplements Linux's SYN-cookie
mechanism (\texttt{net/ipv4/syncookies.c}): when the backlog would overflow, the
server stores \emph{no} half-open state and instead encodes the connection's
identity, a coarse timestamp, and the client's MSS into the 32-bit initial
sequence number of the \texttt{SYN-ACK} (the ``cookie''). Only a client that
truly received the \texttt{SYN-ACK} can echo \texttt{cookie+1} in its
\texttt{ACK}, so the server recomputes and verifies the cookie and rebuilds the
connection --- while a spoofed \texttt{SYN} costs one \texttt{SYN-ACK} and zero
memory.

\begin{lstlisting}[language=bash]
# Enable the real kernel defense that protects nginx (the "protected" trial):
sudo sysctl -w net.ipv4.tcp_syncookies=1
# Demonstrate the project's own reimplementation answering a flood statelessly:
python3 -m server.syn_cookies --connections 1000
\end{lstlisting}

\screenshot{syncookies_demo.png}{\texttt{server/syn\_cookies.py} answering 1000 spoofed SYNs: all legitimate ACKs verified, forged/expired cookies rejected, half-open entries stored = 0.}{cookiedemo}
\screenshot{syncookies_nstat.png}{Victim (protected): \texttt{nstat} showing \texttt{TcpExtSyncookiesSent} rising while HTTP service stays healthy.}{nstat}

\paragraph{Why cookies avoid the flaw of Countermeasure 1.} SYN cookies never
block anyone. They do not decide who is an attacker; they simply refuse to keep
state until a handshake is proven complete. A legitimate client's handshake
succeeds normally, so there is no collateral damage --- which is precisely why
this (and not per-source blocking) is the primary, real-world defense.

\subsection{Comparison}
\begin{table}[H]
\centering
\caption{The two countermeasures side by side.}
\begin{tabular}{@{}lll@{}}
\toprule
 & \textbf{\texttt{defense.py} (SYN Guard)} & \textbf{\texttt{syn\_cookies.py} / kernel} \\ \midrule
Strategy        & detect \& block noisy source IPs   & store no half-open state \\
State kept      & one timer per blocked IP           & none \\
Blocks legit users? & \textcolor{red}{Yes (side effect)} & No \\
Best against    & a few real, high-volume sources    & spoofed-source floods \\
Kernel analogue & none (application level)           & \texttt{net.ipv4.tcp\_syncookies} \\ \bottomrule
\end{tabular}
\end{table}

% ===========================================================================
\section{Results Summary}
% ===========================================================================
Across the three trials --- \textbf{baseline}, \textbf{unprotected} (attack, no
defense), and \textbf{protected} (attack with SYN cookies) --- we compare the
victim's \texttt{syn\_recv\_count} and the legitimate probe's success/latency
(Table~\ref{tab:http}). The expected and observed story is: the flood saturates
the half-open queue and degrades HTTP when undefended; with SYN cookies the
flood still arrives (\texttt{SyncookiesSent} rises) but HTTP service is
preserved; and \texttt{defense.py} suppresses a concentrated source at the cost
of potentially blocking legitimate addresses.

% ===========================================================================
\section{Conclusion}
% ===========================================================================
We reproduced a TCP SYN-flood attack end to end on our own isolated network,
measured its effect on both the kernel's half-open queue and real HTTP service,
and defended against it two ways. The per-source blocker
(\texttt{defense.py}) is observable and effective against concentrated sources
but is fundamentally limited: because it acts on the spoofable source-IP field,
it can block legitimate users. Stateless SYN cookies
(\texttt{syn\_cookies.py} / \texttt{net.ipv4.tcp\_syncookies}) remove the
targeted resource entirely and defend without that collateral damage, which is
why they are the primary real-world countermeasure.

% ===========================================================================
\appendix
\section{Reproduction Commands}
% ===========================================================================
\begin{lstlisting}[language=bash]
# --- Attacker (10.42.0.157) ---
sudo python3 -m client.generator --continuous --rate 1000 --workers 8

# --- Server / victim (10.42.0.81) ---
python3 -m server.monitor --interval 0.5 --output results/attack-monitor.csv
watch -n0.5 'ss -Hnt state syn-recv "( sport = :80 )" | wc -l'
sudo tcpdump -ni <iface> 'tcp port 80 and tcp[tcpflags] & tcp-syn != 0'

# --- Legitimate user probe (run on the server, against its own IP) ---
python3 -m client.sustained_http_load --url http://10.42.0.81/ \
        --duration 20 --timeout 2 --concurrency 10 --output results/attack-http.csv

# --- Defense 1: per-source blocker ---
sudo python3 -m server.defense --interval 0.5 --threshold 5 \
     --block-seconds 30 --output results/defense-events.csv

# --- Defense 2: SYN cookies ---
sudo sysctl -w net.ipv4.tcp_syncookies=1          # protect nginx
python3 -m server.syn_cookies --connections 1000  # demonstrate the mechanism
\end{lstlisting}

\end{document}
